%% kapa: exam=1 round=24 year=2013 period=2 form=A subject=appraisal_law subject_ko=부동산관계법규 seq=39 number=79
%% source: Q-Net 감정평가사 기출문제 (cst00309) artlSeq=5073802 file=extracted/1008/5073802/2013년 제24회 감정평가사 1차시험 기출문제/2교시(A형).hwp
%% conversion: kordoc 4.24.51 -> tools/gen1.py (auto); review=proofread-sonnet; flags=none
\begin{question}{79}
부동산등기법상 등기의 단독신청에 관한 설명으로 옳지 않은 것은?
\begin{choices}
  \item 등기명의인표시의 변경이나 경정의 등기는 해당 권리의 등기명의인이 단독으로 신청한다.
  \item 소유권보존등기 또는 그 말소등기는 등기명의인으로 될 자 또는 등기명의인이 단독으로 신청한다.
  \item 가등기명의인은 단독으로 가등기의 말소를 신청할 수 없다.
  \item 판결에 의한 등기는 승소한 등기권리자 또는 등기의무자가 단독으로 신청한다.
  \item 부동산표시의 변경이나 경정의 등기는 소유권의 등기명의인이 단독으로 신청한다.
\end{choices}
\end{question}
\answer{3}
